\section{AHSME 1983}
        \begin{enumerate}
        	\item (\textit{AHSME 1983}) If $x \neq 0, \frac x{2} = y^2$ and $\frac{x}{4} = 4y$, then $x$ equals


 $\textbf{(A)}\ 8\qquad \textbf{(B)}\ 16\qquad \textbf{(C)}\ 32\qquad \textbf{(D)}\ 64\qquad \textbf{(E)}\ 128$


	\item (\textit{AHSME 1983}) Point $P$ is outside circle $C$ on the plane. At most how many points on $C$ are $3$ cm from $P$? 


 $\textbf{(A)} \ 1 \qquad  \textbf{(B)} \ 2 \qquad  \textbf{(C)} \ 3 \qquad  \textbf{(D)} \ 4 \qquad  \textbf{(E)} \ 8$


	\item (\textit{AHSME 1983}) Three primes $p,q$ and $r$ satisfy $p+q = r$ and $1 < p < q$. Then $p$ equals


 $\textbf{(A)}\ 2\qquad \textbf{(B)}\ 3\qquad \textbf{(C)}\ 7\qquad \textbf{(D)}\ 13\qquad \textbf{(E)}\ 17$


	\item (\textit{AHSME 1983}) 
                    \begin{figure}[H]
                        \centering
                        \includegraphics[width=0.5\textwidth]{Problem/AHSME_I_1983_4_Problem_Figure_1.png}
                    \end{figure}
                    


 In the adjoining plane figure, sides $AF$ and $CD$ are parallel, as are sides $AB$ and $EF$, 
and sides $BC$ and $ED$. Each side has length $1$. Also, $\angle FAB = \angle BCD = 60^\circ$. 
The area of the figure is 


 $\textbf{(A)} \ \frac{\sqrt 3}{2} \qquad  \textbf{(B)} \ 1 \qquad  \textbf{(C)} \ \frac{3}{2} \qquad  \textbf{(D)}\ \sqrt{3}\qquad \textbf{(E)}\ 2$


	\item (\textit{AHSME 1983}) Triangle $ABC$ has a right angle at $C$. If $\sin A = \frac{2}{3}$, then $\tan B$ is


 $\textbf{(A)}\ \frac{3}{5}\qquad \textbf{(B)}\ \frac{\sqrt 5}{3}\qquad \textbf{(C)}\ \frac{2}{\sqrt 5}\qquad \textbf{(D)}\ \frac{\sqrt{5}}{2}\qquad \textbf{(E)}\ \frac{5}{3}$


	\item (\textit{AHSME 1983}) When $x^5, x+\frac{1}{x}$ and $1+\frac{2}{x} + \frac{3}{x^2}$ are multiplied, the product is a polynomial of degree


 $\textbf{(A)}\ 2\qquad \textbf{(B)}\ 3\qquad \textbf{(C)}\ 6\qquad \textbf{(D)}\ 7\qquad \textbf{(E)}\ 8$


	\item (\textit{AHSME 1983}) Alice sells an item at $\$10$ less than the list price and receives $10\%$ of her selling price as her commission. 
Bob sells the same item at $\$20$ less than the list price and receives $20\%$ of his selling price as his commission. 
If they both get the same commission, then the list price is


 $\textbf{(A) } \$20\qquad \textbf{(B) } \$30\qquad \textbf{(C) } \$50\qquad \textbf{(D) } \$70\qquad \textbf{(E) } \$100$


	\item (\textit{AHSME 1983}) Let $f(x) = \frac{x+1}{x-1}$. Then for $x^2 \neq 1, f(-x)$ is


 $\textbf{(A)}\ \frac{1}{f(x)}\qquad \textbf{(B)}\ -f(x)\qquad \textbf{(C)}\ \frac{1}{f(-x)}\qquad \textbf{(D)}\ -f(-x)\qquad \textbf{(E)}\ f(x)$


	\item (\textit{AHSME 1983}) In a certain population the ratio of the number of women to the number of men is $11$ to $10$. 
If the average (arithmetic mean) age of the women is $34$ and the average age of the men is $32$, 
then the average age of the population is


 $\textbf{(A)}\ 32\frac{9}{10}\qquad \textbf{(B)}\ 32\frac{20}{21}\qquad \textbf{(C)}\ 33\qquad \textbf{(D)}\ 33\frac{1}{21}\qquad \textbf{(E)}\ 33\frac{1}{10}$


	\item (\textit{AHSME 1983}) Segment $AB$ is both a diameter of a circle of radius $1$ and a side of an equilateral triangle $ABC$. 
The circle also intersects $AC$ and $BC$ at points $D$ and $E$, respectively. The length of $AE$ is 


 $\textbf{(A)} \ \frac{3}{2} \qquad  \textbf{(B)} \ \frac{5}{3} \qquad  \textbf{(C)} \ \frac{\sqrt 3}{2} \qquad  \textbf{(D)}\ \sqrt{3}\qquad \textbf{(E)}\ \frac{2+\sqrt 3}{2}$


	\item (\textit{AHSME 1983}) Simplify $\sin (x-y) \cos y + \cos (x-y) \sin y$.


 $\textbf{(A)}\ 1\qquad \textbf{(B)}\ \sin x\qquad \textbf{(C)}\ \cos x\qquad \textbf{(D)}\ \sin x \cos 2y\qquad \textbf{(E)}\ \cos x\cos 2y$


	\item (\textit{AHSME 1983}) If $\log_7 \Big(\log_3 (\log_2 x) \Big) = 0$, then $x^{-1/2}$ equals 


 $\textbf{(A)} \ \frac{1}{3} \qquad  \textbf{(B)} \ \frac{1}{2 \sqrt 3} \qquad  \textbf{(C)}\ \frac{1}{3\sqrt 3}\qquad \textbf{(D)}\ \frac{1}{\sqrt{42}}\qquad \textbf{(E)}\ \text{none of these}$


	\item (\textit{AHSME 1983}) If $xy = a, xz =b,$ and $yz = c$, and none of these quantities is $0$, then $x^2+y^2+z^2$ equals


 $\textbf{(A)}\ \frac{ab+ac+bc}{abc}\qquad \textbf{(B)}\ \frac{a^2+b^2+c^2}{abc}\qquad \textbf{(C)}\ \frac{(a+b+c)^2}{abc}\qquad \textbf{(D)}\ \frac{(ab+ac+bc)^2}{abc}\qquad \textbf{(E)}\ \frac{(ab)^2+(ac)^2+(bc)^2}{abc}$


	\item (\textit{AHSME 1983}) The units digit of $3^{1001} 7^{1002} 13^{1003}$ is


 $\textbf{(A)}\ 1\qquad \textbf{(B)}\ 3\qquad \textbf{(C)}\ 5\qquad \textbf{(D)}\ 7\qquad \textbf{(E)}\ 9$


	\item (\textit{AHSME 1983}) Three balls marked $1,2$ and $3$ are placed in an urn. One ball is drawn, its number is recorded, and then the ball is returned to the urn. This process is repeated and then repeated once more, and each ball is equally likely to be drawn on each occasion. If the sum of the numbers recorded is $6$, what is the probability that the ball numbered $2$ was drawn all three times? 


 $\textbf{(A)} \ \frac{1}{27} \qquad  \textbf{(B)} \ \frac{1}{8} \qquad  \textbf{(C)} \ \frac{1}{7} \qquad  \textbf{(D)} \ \frac{1}{6} \qquad \textbf{(E)}\ \frac{1}{3}$


	\item (\textit{AHSME 1983}) Let $x = .123456789101112....998999$, where the digits are obtained by writing the integers $1$ through $999$ in order. 
The $1983$<sup>rd</sup> digit to the right of the decimal point is


 $\textbf{(A)}\ 2\qquad \textbf{(B)}\ 3\qquad \textbf{(C)}\ 5\qquad \textbf{(D)}\ 7\qquad \textbf{(E)}\ 8$


	\item (\textit{AHSME 1983}) 
                    \begin{figure}[H]
                        \centering
                        \includegraphics[width=0.5\textwidth]{Problem/AHSME_I_1983_17_Problem_Figure_1.png}
                    \end{figure}
                    


 The diagram above shows several numbers in the complex plane. The circle is the unit circle centered at the origin. 
One of these numbers is the reciprocal of $F$. Which one? 


 $\textbf{(A)} \ A \qquad  \textbf{(B)} \ B \qquad  \textbf{(C)} \ C \qquad  \textbf{(D)} \ D \qquad  \textbf{(E)} \ E$


	\item (\textit{AHSME 1983}) Let $f$ be a polynomial function such that, for all real $x$,
$f(x^2 + 1) = x^4 + 5x^2 + 3$.
For all real $x, f(x^2-1)$ is


 $\textbf{(A)}\ x^4+5x^2+1\qquad \textbf{(B)}\ x^4+x^2-3\qquad \textbf{(C)}\ x^4-5x^2+1\qquad \textbf{(D)}\ x^4+x^2+3\qquad \textbf{(E)}\ \text{none of these}$


	\item (\textit{AHSME 1983}) Point $D$ is on side $CB$ of triangle $ABC$. If 
$\angle{CAD} = \angle{DAB} = 60^\circ\mbox{, }AC = 3\mbox{ and }AB = 6$, 
then the length of $AD$ is 


 $\textbf{(A)} \ 2 \qquad  \textbf{(B)} \ 2.5 \qquad  \textbf{(C)} \ 3 \qquad  \textbf{(D)} \ 3.5 \qquad  \textbf{(E)} \ 4$


	\item (\textit{AHSME 1983}) If $\tan{\alpha}$ and $\tan{\beta}$ are the roots of $x^2 - px + q = 0$, and $\cot{\alpha}$ and $\cot{\beta}$
are the roots of $x^2 - rx + s = 0$, then $rs$ is necessarily 


 $\textbf{(A)} \ pq \qquad  \textbf{(B)} \ \frac{1}{pq} \qquad  \textbf{(C)} \ \frac{p}{q^2} \qquad  \textbf{(D)}\ \frac{q}{p^2}\qquad \textbf{(E)}\ \frac{p}{q}$


	\item (\textit{AHSME 1983}) Find the smallest positive number from the numbers below.


 $\textbf{(A)} \ 10-3\sqrt{11} \qquad  \textbf{(B)} \ 3\sqrt{11}-10 \qquad  \textbf{(C)}\ 18-5\sqrt{13}\qquad \textbf{(D)}\ 51-10\sqrt{26}\qquad \textbf{(E)}\ 10\sqrt{26}-51$


	\item (\textit{AHSME 1983}) Consider the two functions $f(x) = x^2+2bx+1$ and $g(x) = 2a(x+b)$, where the variable $x$ and the constants $a$ and $b$ are real numbers. 
Each such pair of constants $a$ and $b$ may be considered as a point $(a,b)$ in an $ab$-plane. 
Let $S$ be the set of such points $(a,b)$ for which the graphs of $y = f(x)$ and $y = g(x)$ do \textbf{not} intersect (in the $xy$-plane). The area of $S$ is 


 $\textbf{(A)} \ 1 \qquad  \textbf{(B)} \ \pi \qquad  \textbf{(C)} \ 4 \qquad  \textbf{(D)} \ 4 \pi \qquad  \textbf{(E)} \ \text{infinite}$


	\item (\textit{AHSME 1983}) In the adjoining figure the five circles are tangent to one another consecutively and to the lines 
$L_1$ and $L_2$. 
If the radius of the largest circle is $18$ and that of the smallest one is $8$, then the radius of the middle circle is


 
                    \begin{figure}[H]
                        \centering
                        \includegraphics[width=0.5\textwidth]{Problem/AHSME_I_1983_23_Problem_Figure_1.png}
                    \end{figure}
                    
$\textbf{(A)} \ 12 \qquad  \textbf{(B)} \ 12.5 \qquad  \textbf{(C)} \ 13 \qquad  \textbf{(D)} \ 13.5 \qquad  \textbf{(E)} \ 14$


	\item (\textit{AHSME 1983}) How many non-congruent right triangles are there such that the perimeter in $\text{cm}$ and the area in $\text{cm}^2$ are numerically equal? 


 $\textbf{(A)} \ \text{none} \qquad  \textbf{(B)} \ 1 \qquad  \textbf{(C)} \ 2 \qquad  \textbf{(D)} \ 4 \qquad  \textbf{(E)} \ \text{infinitely many}$


	\item (\textit{AHSME 1983}) If $60^a = 3$ and $60^b = 5$, then $12^{(1-a-b)/\left(2\left(1-b\right)\right)}$ is


 $\textbf{(A)} \ \sqrt{3} \qquad  \textbf{(B)} \ 2 \qquad  \textbf{(C)} \ \sqrt{5} \qquad  \textbf{(D)} \ 3 \qquad  \textbf{(E)} \ \sqrt{12}$


	\item (\textit{AHSME 1983}) The probability that event $A$ occurs is $\frac{3}{4}$; the probability that event $B$ occurs is $\frac{2}{3}$. 
Let $p$ be the probability that both $A$ and $B$ occur. The smallest interval necessarily containing $p$ is the interval


 $\textbf{(A)}\ \Big[\frac{1}{12},\frac{1}{2}\Big]\qquad \textbf{(B)}\ \Big[\frac{5}{12},\frac{1}{2}\Big]\qquad \textbf{(C)}\ \Big[\frac{1}{2},\frac{2}{3}\Big]\qquad \textbf{(D)}\ \Big[\frac{5}{12},\frac{2}{3}\Big]\qquad \textbf{(E)}\ \Big[\frac{1}{12},\frac{2}{3}\Big]$


	\item (\textit{AHSME 1983}) A large sphere is on a horizontal field on a sunny day. At a certain time the shadow of the sphere reaches out a distance 
of $10$ m from the point where the sphere touches the ground. At the same instant a meter stick 
(held vertically with one end on the ground) casts a shadow of length $2$ m. What is the radius of the sphere in meters? 
(Assume the sun's rays are parallel and the meter stick is a line segment.)


 $\textbf{(A)}\ \frac{5}{2}\qquad \textbf{(B)}\ 9 - 4\sqrt{5}\qquad \textbf{(C)}\ 8\sqrt{10} - 23\qquad \textbf{(D)}\ 6-\sqrt{15}\qquad \textbf{(E)}\ 10\sqrt{5}-20$


	\item (\textit{AHSME 1983}) Triangle $\triangle ABC$ in the figure has area $10$. Points $D, E$ and $F$, all distinct from $A, B$ and $C$, 
are on sides $AB, BC$ and $CA$ respectively, and $AD = 2, DB = 3$. If triangle $\triangle ABE$ and quadrilateral $DBEF$ 
have equal areas, then that area is


 
                    \begin{figure}[H]
                        \centering
                        \includegraphics[width=0.5\textwidth]{Problem/AHSME_I_1983_28_Problem_Figure_1.png}
                    \end{figure}
                    
$\textbf{(A)}\ 4\qquad \textbf{(B)}\ 5\qquad \textbf{(C)}\ 6\qquad \textbf{(D)}\ \frac{5}{3}\sqrt{10}\qquad \textbf{(E)}\ \text{not uniquely determined}$


	\item (\textit{AHSME 1983}) A point $P$ lies in the same plane as a given square of side $1$. Let the vertices of the square, 
taken counterclockwise, be $A, B, C$ and $D$. Also, let the distances from $P$ to $A, B$ and $C$, respectively, be $u, v$ and $w$. 
What is the greatest distance that $P$ can be from $D$ if $u^2 + v^2 = w^2$?


 $\textbf{(A)}\ 1 + \sqrt{2}\qquad \textbf{(B)}\ 2\sqrt{2}\qquad \textbf{(C)}\ 2 + \sqrt{2}\qquad \textbf{(D)}\ 3\sqrt{2}\qquad \textbf{(E)}\ 3+\sqrt{2}$


	\item (\textit{AHSME 1983}) Distinct points $A$ and $B$ are on a semicircle with diameter $MN$ and center $C$. 
The point $P$ is on $CN$ and $\angle CAP = \angle CBP = 10^{\circ}$. If $\stackrel{\frown}{MA} = 40^{\circ}$, then $\stackrel{\frown}{BN}$ equals


 
                    \begin{figure}[H]
                        \centering
                        \includegraphics[width=0.5\textwidth]{Problem/AHSME_I_1983_30_Problem_Figure_1.png}
                    \end{figure}
                    
$\textbf{(A)}\ 10^{\circ}\qquad \textbf{(B)}\ 15^{\circ}\qquad \textbf{(C)}\ 20^{\circ}\qquad \textbf{(D)}\ 25^{\circ}\qquad \textbf{(E)}\ 30^{\circ}$


\end{enumerate}